% ============================================================
%  C: Como Programar -- 6ª Edição | Deitel & Deitel
%  Capítulo 4 -- Controle de Programa em C
%  FAZENDO A DIFERENÇA (4.38, 4.39)
% ============================================================
\documentclass[12pt,a4paper]{article}
\usepackage[utf8]{inputenc}
\usepackage[T1]{fontenc}
\usepackage[brazil]{babel}
\usepackage{geometry}
\geometry{top=2.5cm,bottom=2.5cm,left=3cm,right=3cm}
\usepackage{parskip}\setlength{\parskip}{6pt}\setlength{\parindent}{0pt}
\usepackage{xcolor,tcolorbox,enumitem,listings,fancyhdr,titlesec,hyperref,booktabs,array,amsmath}
\tcbuselibrary{skins,breakable}

\definecolor{deitelblue}{RGB}{0,70,127}
\definecolor{deitelgray}{RGB}{240,242,245}
\definecolor{deitelgreen}{RGB}{0,110,60}
\definecolor{deitellightblue}{RGB}{220,235,250}
\definecolor{deitelred}{RGB}{160,20,20}
\definecolor{deitellorange}{RGB}{200,90,0}
\definecolor{ecogreen}{RGB}{0,120,50}
\definecolor{ecolightgreen}{RGB}{210,240,220}

\newtcolorbox{boaprática}[1][]{colback=deitellightblue,colframe=deitelblue,
  fonttitle=\bfseries\small,title={Boa prática de programação},breakable,#1}
\newtcolorbox{respostabox}[1][]{colback=deitelgray,colframe=deitelblue!60,
  fonttitle=\bfseries\small\color{deitelblue},title={Resposta detalhada},breakable,#1}
\newtcolorbox{enunciadoBox}[1][]{colback=white,colframe=deitelblue,
  fonttitle=\bfseries,title={#1},breakable}
\newtcolorbox{saida}[1][]{colback=black!88,colframe=deitelblue,
  fontupper=\ttfamily\small\color{green!70},
  title={\small\color{white}Saída do programa},breakable,#1}
\newtcolorbox{pesquisa}[1][]{colback=ecolightgreen,colframe=ecogreen,
  fonttitle=\bfseries\small,title={Pesquisa externa realizada},breakable,#1}
\newtcolorbox{atencao}[1][]{colback=yellow!10,colframe=deitellorange,
  fonttitle=\bfseries\small,title={Atenção},breakable,#1}

\setlist[enumerate,1]{label=\textbf{\alph*)},leftmargin=2em}
\setlist[itemize]{leftmargin=2em}

\lstset{
  basicstyle=\ttfamily\small,backgroundcolor=\color{deitelgray},
  frame=single,framerule=0.4pt,rulecolor=\color{deitelblue!40},
  breaklines=true,showstringspaces=false,
  keywordstyle=\color{deitelblue}\bfseries,
  commentstyle=\color{deitelgreen}\itshape,
  stringstyle=\color{deitelred},
  language=C,numbers=left,numberstyle=\tiny\color{gray},
  xleftmargin=18pt,xrightmargin=5pt,stepnumber=1,
}

\pagestyle{fancy}\fancyhf{}
\fancyhead[L]{\small\color{deitelblue}\textbf{C: Como Programar -- 6ª ed. | Deitel \& Deitel}}
\fancyhead[R]{\small\color{deitelblue}Cap. 4 -- Fazendo a Diferença}
\fancyfoot[C]{\small\thepage}
\renewcommand{\headrulewidth}{0.4pt}

\titleformat{\section}[block]{\large\bfseries\color{deitelblue}}{\thesection.}{0.5em}{}[\titlerule]
\titleformat{\subsection}[block]{\normalsize\bfseries\color{deitelblue!80}}{}{}{}

\hypersetup{colorlinks=true,linkcolor=deitelblue,urlcolor=ecogreen}

\begin{document}

\begin{titlepage}
  \centering\vspace*{2cm}
  {\color{deitelblue}\rule{\linewidth}{2pt}}\\[0.4cm]
  {\huge\bfseries\color{deitelblue}C: Como Programar}\\[0.2cm]
  {\large\color{deitelblue}6ª Edição $\cdot$ Paul Deitel \& Harvey Deitel}\\[0.2cm]
  {\color{deitelblue}\rule{\linewidth}{2pt}}\\[1cm]
  {\LARGE\bfseries Capítulo 4}\\[0.3cm]
  {\Large\color{deitelblue!80}Controle de Programa em C}\\[0.8cm]
  \begin{tcolorbox}[colback=ecolightgreen,colframe=ecogreen,width=0.85\linewidth]
    \centering{\large\bfseries\color{ecogreen}Fazendo a Diferença}\\[0.3cm]
    {\normalsize Exercícios 4.38 e 4.39\\[0.2cm]
    Crescimento populacional ao longo de 75 anos\\
    Análise da proposta tributária ``FairTax''}
  \end{tcolorbox}
  \vfill{\small Abordagem incremental Deitel}
\end{titlepage}

\tableofcontents\newpage

% ============================================================
\section{Exercício 4.38 -- Crescimento Populacional Mundial em 75 Anos}
% ============================================================

\begin{enunciadoBox}[Enunciado 4.38]
Pesquise sobre crescimento populacional. Encontre a população mundial atual e sua taxa
anual de crescimento. Escreva um programa que projete o crescimento populacional pelos
próximos 75 anos, exibindo: (1) ano, (2) população projetada ao final desse ano, e
(3) aumento numérico ocorrido nesse ano. Determine em que ano a população dobraria,
considerando taxa constante.
\end{enunciadoBox}

\subsection*{Pesquisa Externa}

\begin{pesquisa}
Conforme dados das Nações Unidas (UN/DESA), a população mundial atual em 2024-2025
era de aproximadamente \textbf{8,1 bilhões de pessoas}. A taxa de crescimento anual
médio nos últimos anos foi de cerca de \textbf{0,87\%} -- uma desaceleração significativa
em relação aos picos de mais de 2\% ao ano nas décadas de 1960-70.

\textbf{Tendência futura:} Projeções da ONU sugerem que a taxa continuará caindo, e
a população mundial pode atingir um pico de \(\sim 10{,}4\) bilhões em torno de 2086,
para então começar a declinar. Este exercício, no entanto, pede para usar uma taxa
\textit{constante}, criando um cenário hipotético de comparação.

\textbf{Fórmula de crescimento composto:}
\[
  P_n = P_{n-1} \cdot (1 + r)
\]
ou, de forma fechada, $P_n = P_0 \cdot (1+r)^n$.

\textbf{Tempo de duplicação} (regra dos 70):
\[
  t_{\text{dup}} \approx \frac{\ln 2}{\ln(1+r)} \approx \frac{70}{r \cdot 100}
\]
Para $r = 0{,}87\%$: $t_{\text{dup}} \approx 70 / 0{,}87 \approx 80{,}5$ anos.
\end{pesquisa}

\subsection*{Solução em C}

\begin{respostabox}
\begin{lstlisting}
/* Ex4.38: crescimento_populacional_75anos.c
   Tabela de crescimento populacional para os proximos 75 anos */
#include <stdio.h>

int main( void )
{
    double populacao;          /* populacao em bilhoes      */
    double taxa;               /* taxa anual em decimal     */
    double populacaoInicial;   /* guarda valor original     */
    double aumento;            /* aumento absoluto no ano   */
    double populacaoAnterior;  /* para calcular o aumento   */
    int ano;
    int anoDuplicacao = 0;     /* ano em que dobra (0 = ainda nao) */

    printf( "Digite a populacao mundial atual (em bilhoes): " );
    scanf( "%lf", &populacao );

    printf( "Digite a taxa anual de crescimento (decimal, "
            "ex: 0.0087 para 0.87%%): " );
    scanf( "%lf", &taxa );

    populacaoInicial = populacao;

    /* Cabecalho da tabela */
    printf( "\n%5s %20s %20s\n", "Ano", "Populacao (bilhoes)",
            "Aumento (bilhoes)" );
    printf( "----- -------------------- --------------------\n" );

    /* Projecao para 75 anos */
    for ( ano = 1; ano <= 75; ++ano ) {
        populacaoAnterior = populacao;
        populacao = populacao * ( 1.0 + taxa );
        aumento = populacao - populacaoAnterior;

        printf( "%5d %20.4f %20.4f\n", ano, populacao, aumento );

        /* Verifica primeiro ano em que a populacao dobra */
        if ( anoDuplicacao == 0 && populacao >= 2.0 * populacaoInicial ) {
            anoDuplicacao = ano;
        }
    }

    if ( anoDuplicacao > 0 ) {
        printf( "\nA populacao DOBROU no ano %d "
                "(%.4f bilhoes).\n", anoDuplicacao,
                populacaoInicial * 2.0 );
    }
    else {
        printf( "\nA populacao NAO dobrou nos proximos 75 anos.\n" );
    }

    return 0;
} /* fim da funcao main */
\end{lstlisting}

\textbf{Trecho da saída esperada} (com $P_0 = 8{,}1$ bilhões e $r = 0{,}87\%$):

\begin{saida}
\hspace{0.3em}Ano \hspace{1.5em}Populacao (bilhoes) \hspace{1em}Aumento (bilhoes)\\
----- -------------------- --------------------\\
\hspace{2em}1 \hspace{8em}8.1705 \hspace{8em}0.0705\\
\hspace{2em}2 \hspace{8em}8.2416 \hspace{8em}0.0711\\
\hspace{2em}3 \hspace{8em}8.3133 \hspace{8em}0.0717\\
\ldots\\
\hspace{1.5em}75 \hspace{7.5em}15.5128 \hspace{8em}0.1339\\
\\
A populacao DOBROU no ano 80 (16.2000 bilhoes).
\end{saida}

\textit{Observação:} Com a taxa de 0,87\% ao ano, a duplicação ocorreria por volta
do ano \textbf{80} (após 75 anos). Para taxas mais altas, dobrar leva menos tempo
(regra dos 70 / taxa\%).
\end{respostabox}

\subsection*{Reflexão Crítica}

\begin{tcolorbox}[colback=ecolightgreen,colframe=ecogreen,
  title={\bfseries Por que esse exercício importa?}]

A simulação mostra um fato crucial sobre crescimento exponencial: \textit{mesmo
taxas pequenas geram aumentos enormes ao longo do tempo}. Mais importante ainda,
o \textbf{aumento absoluto ano a ano cresce}: começamos adicionando 70 milhões
de pessoas no ano 1 e terminamos adicionando 134 milhões no ano 75!

Isso suscita reflexões importantes:
\begin{itemize}[noitemsep]
  \item \textbf{Recursos finitos:} terra arável, água potável, ar respirável -- todos
  têm limites físicos. Uma população crescendo exponencialmente em um planeta finito
  enfrenta tensões inevitáveis.
  \item \textbf{Sustentabilidade:} este exercício se conecta diretamente aos exercícios
  de pegada de carbono (1.10) e transporte solidário (2.33). Mais pessoas significa
  mais emissões -- a menos que mudemos nossos padrões de consumo.
  \item \textbf{Realidade vs.\ modelo:} a hipótese de \textit{taxa constante} é uma
  simplificação. Em escolas avançadas de demografia, modelos logísticos
  (curvas em S) capturam melhor o comportamento real, com saturação e declínio.
\end{itemize}

\textbf{Conexão com programação:} este programa é um exemplo perfeito de uso da
estrutura \texttt{for} para simulações temporais -- um padrão recorrente em finanças,
biologia, física e engenharia.
\end{tcolorbox}

\newpage

% ============================================================
\section{Exercício 4.39 -- Análise da Proposta ``FairTax''}
% ============================================================

\begin{enunciadoBox}[Enunciado 4.39]
A proposta \textbf{FairTax}, nos EUA, sugere eliminar o imposto de renda em favor
de uma taxa de consumo de 23\% sobre todos os bens e serviços. Críticos argumentam
que, devido ao modo de cálculo, a alíquota efetiva seria de 30\%. Investigue ambas as
visões e escreva um programa que peça os gastos do usuário em diversas categorias
e calcule o imposto FairTax estimado.
\end{enunciadoBox}

\subsection*{Pesquisa: Os Dois Métodos de Cálculo}

\begin{pesquisa}
A controvérsia sobre se a FairTax é ``23\%'' ou ``30\%'' decorre de duas formas
distintas de calcular percentuais sobre vendas:

\textbf{Método 1 -- Taxa inclusiva (\textit{tax-inclusive}):}\\
A alíquota é expressa como percentual do preço \textbf{final pago} (que já inclui
o imposto). Exemplo: para um produto de \$100, com taxa inclusiva de 23\%:
\[
  \text{imposto} = 0{,}23 \times \$100 = \$23
\]
O comprador paga \$100, dos quais \$77 vão ao vendedor e \$23 ao governo.

\textbf{Método 2 -- Taxa exclusiva (\textit{tax-exclusive}):}\\
A alíquota é expressa como percentual do preço \textbf{antes do imposto} (preço base
que o vendedor recebe). É a forma usual de impostos sobre vendas nos EUA. Para receber
\$77 e pagar \$23 de imposto, a alíquota exclusiva é:
\[
  \text{taxa\_exclusiva} = \frac{23}{77} \approx 29{,}87\%
\]

\textbf{Conclusão:} ambas as taxas representam o \textit{mesmo} cenário -- cobrança de
\$23 sobre uma transação total de \$100 (ou \$23 sobre um preço base de \$77). A diferença
é apenas a base de cálculo. Apoiadores da FairTax preferem ``23\%'' (inclusiva) porque
parece menor; críticos preferem ``30\%'' (exclusiva) porque ela é mais comparável às
alíquotas tradicionais de impostos sobre vendas estaduais nos EUA.

Para gastos totais $G$, o imposto pago (taxa inclusiva 23\%) é:
\[
  \text{imposto} = 0{,}23 \cdot G
\]
\end{pesquisa}

\subsection*{Solução em C}

\begin{respostabox}
\begin{lstlisting}
/* Ex4.39: fairtax.c
   Calcula imposto FairTax estimado a partir de gastos por categoria */
#include <stdio.h>

int main( void )
{
    int categoria;
    float gasto;
    float gastosCategorias[ 8 ] = { 0 }; /* nao usaremos - so soma */
    float total = 0.0f;
    float impostoInclusivo;   /* taxa de 23% */
    float impostoExclusivo;   /* taxa de ~30% sobre o preco base */
    float gastoBase;          /* gasto sem o imposto */

    /* Categorias predefinidas */
    const int NUM_CATEGORIAS = 7;

    printf( "===== Calculadora FairTax =====\n\n" );
    printf( "Informe seus gastos anuais nas categorias abaixo.\n" );
    printf( "(Os valores devem ser totais ja com tributos atuais)\n\n" );

    for ( categoria = 1; categoria <= NUM_CATEGORIAS; ++categoria ) {
        switch ( categoria ) {
            case 1: printf( "Moradia (R$): " );      break;
            case 2: printf( "Alimentacao (R$): " );  break;
            case 3: printf( "Vestuario (R$): " );    break;
            case 4: printf( "Transporte (R$): " );   break;
            case 5: printf( "Educacao (R$): " );     break;
            case 6: printf( "Saude (R$): " );        break;
            case 7: printf( "Ferias (R$): " );       break;
        }
        scanf( "%f", &gasto );
        total += gasto;
    }

    /* Calculo da FairTax pelos dois metodos */
    impostoInclusivo = total * 0.23f;            /* 23% sobre o total */
    gastoBase        = total - impostoInclusivo;  /* preco base */
    impostoExclusivo = total * ( 23.0f / 77.0f ); /* equivale a 29,87% */

    /* Note: ambas as formulas dao o MESMO valor monetario,
       sao apenas DUAS FORMAS de descrever a mesma alíquota */

    printf( "\n========== RELATORIO ==========\n" );
    printf( "Total de gastos:        R$ %10.2f\n", total );
    printf( "-------------------------------\n" );
    printf( "Imposto FairTax estimado:\n" );
    printf( "  Taxa inclusiva (23%%): R$ %10.2f\n", impostoInclusivo );
    printf( "  (~30%% taxa exclusiva): R$ %10.2f\n", impostoExclusivo );
    printf( "  Gasto liquido (preco base): R$ %10.2f\n", gastoBase );
    printf( "===============================\n\n" );

    printf( "INTERPRETACAO:\n" );
    printf( "  - 23%% inclusiva: o imposto representa 23%% do total pago.\n" );
    printf( "  - 30%% exclusiva: o imposto representa 30%% do preco base.\n" );
    printf( "  Ambos descrevem o MESMO valor monetario (R$ %.2f).\n",
            impostoInclusivo );

    return 0;
} /* fim da funcao main */
\end{lstlisting}

\begin{saida}
===== Calculadora FairTax =====\\[2pt]
Informe seus gastos anuais nas categorias abaixo.\\[2pt]
Moradia (R\$): 24000\\
Alimentacao (R\$): 12000\\
Vestuario (R\$): 3000\\
Transporte (R\$): 9000\\
Educacao (R\$): 6000\\
Saude (R\$): 5000\\
Ferias (R\$): 4000\\[2pt]
========== RELATORIO ==========\\
Total de gastos: \hspace{4em} R\$ \hspace{0.5em}63000.00\\
-------------------------------\\
Imposto FairTax estimado:\\
\hphantom{xx}Taxa inclusiva (23\%): \hspace{0.4em} R\$ \hspace{1em}14490.00\\
\hphantom{xx}(~30\% taxa exclusiva): R\$ \hspace{1em}18818.18\\
\hphantom{xx}Gasto liquido (preco base): R\$ \hspace{1em}48510.00\\
===============================
\end{saida}
\end{respostabox}

\subsection*{Análise Matemática Detalhada}

\begin{atencao}
\textbf{Atenção à apresentação dos dois métodos:} as duas alíquotas (23\% inclusiva
e 29,87\% exclusiva) descrevem \textit{exatamente o mesmo dinheiro}. A
\textbf{diferença é apenas semântica}, sobre qual é a ``base de comparação''.

Veja com um exemplo concreto:
\begin{itemize}[noitemsep]
  \item Você gasta \textbf{R\$ 100} no total na compra (preço de etiqueta + imposto).
  \item Desses R\$ 100: \textbf{R\$ 23} vão para o governo e \textbf{R\$ 77} para o
  vendedor.
  \item ``23\% inclusivo'': $\frac{23}{100} = 23\%$ -- correto.
  \item ``30\% exclusivo'': $\frac{23}{77} = 29{,}87\%$ -- também correto.
\end{itemize}

\textbf{Por que é importante essa distinção?} Porque a maioria dos impostos sobre
vendas nos EUA (\textit{sales tax} estaduais) é tradicionalmente expressa de forma
\textbf{exclusiva}: você compra um item de \$100 e paga \$108 com 8\% de \textit{sales
tax}. Quando a FairTax usa o método inclusivo, há uma percepção visual de alíquota
menor (23\% < 30\%) -- daí a controvérsia.
\end{atencao}

\subsection*{Reflexão Sobre Sistemas Tributários}

\begin{tcolorbox}[colback=ecolightgreen,colframe=ecogreen,
  title={\bfseries Computação como ferramenta para entendimento cívico}]

Este exercício, aparentemente um simples programa de soma e multiplicação, ilustra
como a programação pode ser uma ferramenta poderosa para \textbf{compreender
políticas públicas}. Ao construir o cálculo, somos forçados a entender:
\begin{itemize}[noitemsep]
  \item Como impostos são expressos matematicamente
  \item Por que a apresentação retórica importa
  \item Como duas afirmações ``ambas verdadeiras'' podem dar impressões opostas
\end{itemize}

A capacidade de programar e analisar quantitativamente políticas é uma das marcas
do cidadão informado no século XXI. Os Deitel chamam isso de ``\textit{making a
difference}'' -- usar a programação não apenas para resolver problemas técnicos,
mas para participar mais conscientemente das discussões da sociedade.
\end{tcolorbox}

\newpage

% ============================================================
\section*{Resumo -- Fazendo a Diferença no Capítulo~4}

\begin{tcolorbox}[colback=deitellightblue,colframe=deitelblue,
  title={\bfseries Conexões e progressão}]

\textbf{4.38 -- Crescimento Populacional Mundial:}
Aplicação de \texttt{for} a uma simulação temporal de longo prazo. Conecta-se diretamente
ao Ex.~3.47 (versão de 5 anos) -- agora estendido para 75 anos com tabela detalhada
e detecção do ano de duplicação. Demonstra o poder do laço \texttt{for} para
modelagem matemática iterativa.

\medskip\textbf{4.39 -- FairTax:}
Uso da estrutura \texttt{switch} (Cap.~4) para tratar múltiplas categorias de gastos
de forma elegante. A questão matemática (taxa inclusiva vs.\ exclusiva) é uma lição
de literacia financeira aplicada a programação.

\medskip\textbf{Padrão recorrente nesta seção:} ambos os exercícios usam
\textit{simulações iterativas} -- a essência do que computadores fazem melhor que
humanos. Calcular 75 iterações de crescimento a mão seria tedioso e propenso a
erros; o computador faz isso em frações de segundo, com perfeita precisão.

\end{tcolorbox}

\begin{boaprática}
  Os exercícios ``Fazendo a Diferença'' do Capítulo~4 demonstram como, com apenas as
  ferramentas dos primeiros 4 capítulos -- variáveis, \texttt{if/else}, \texttt{for/while},
  \texttt{switch} e operações aritméticas -- já podemos abordar problemas de relevância
  social e econômica significativa. Cada novo capítulo expandirá ainda mais o seu
  arsenal de ferramentas para fazer a diferença.
\end{boaprática}

\end{document}
